%% =====================================================================
%%  第 10 章　内耗
%%  源文件：10-内耗.md
%% =====================================================================
\chapter{内耗}

\applicable{一件事已经发生，或者一条消息已经发出，而你还在反复地想它。}

\begin{guideblock}
\textbf{本章解决什么问题}：为什么你的请求早就发出去了，客户端还在跑；以及怎么把这个进程停下来。

先给结论：\textbf{内耗是一次本地重试风暴。}请求已经发出，返回还没到，而你的客户端不肯等——它在本地反复重跑这条请求，用自己造的返回值，替换掉那个还没到的真实返回值。

值得注意的是，本地生成的返回有一个共同特征：\textbf{它永远比真实返回更坏。}这不是因为它更准，而是因为造它的模型，是用你的恐惧训练出来的。

需要说明的是，本章不处理“如何成为一个不焦虑的人”这个问题——那构成了一个无法执行的任务。本章处理的，是一个更小的、可以执行的任务：\textbf{识别这次重试，并把它停掉。}
\end{guideblock}

\hcirule

\section{请求已经发出去了}

先看样本。

\sample{1010}\par
\begin{mdquote}
晚上十一点，你给一个朋友发了条消息：“周末那事，你还去吗？”\par
\vspace{0.3\baselineskip}
他没回。\par
\vspace{0.3\baselineskip}
接下来三个小时，你在脑子里过了一遍：他是不是不高兴了？是不是我上次哪句说错了？他会不会觉得我烦？要不我撤回——撤回更奇怪吧。\par
\vspace{0.3\baselineskip}
凌晨两点，你睡着了。\par
\vspace{0.3\baselineskip}
第二天早上八点，他回：“去啊，昨天睡了，不好意思。”
\end{mdquote}

拆解这次交互。

你的请求在 23:00 发出。对方的真实返回在次日 08:00 到达。\textbf{两者之间隔了九个小时。}

而在那九个小时里，你的本地进程跑了大约几十次重试。

\textbf{问题不在那九个小时。}问题在于，你在这段时间里，用本地生成的返回值替换掉了那个尚未到达的真实返回值——然后\textbf{按照这个假返回值，调整了你的行为}：你可能真的没再主动联系他，可能在他回复后语气变冷，也可能在别处跟人提起这件事。

这里需要给“重试风暴”一句人话解释：\textbf{一次正常的重试，是等一段时间、再发一次请求；一次重试风暴，是不等返回、关掉超时、然后用自己算出来的结果向下继续。}前者改变的是网络上发生的事情，后者改变的是你脑子里发生的事情。

再看一个更平常的样本。

\sample{1011}\par
\begin{mdquote}
周三下午两点，你把一份改好的方案发给了合作方的对接人，附了一句“麻烦看下有没有问题”。\par
\vspace{0.3\baselineskip}
三点，没有回。你想：可能在看。\par
\vspace{0.3\baselineskip}
四点，没有回。你想：是不是我哪里改错了。\par
\vspace{0.3\baselineskip}
五点，你翻了一遍两人的历史聊天记录，找出了三处“他之前也没回”的证据。\par
\vspace{0.3\baselineskip}
六点，你决定明天再说，并顺手把方案的措辞又改了一遍——尽管他一个字都没提。\par
\vspace{0.3\baselineskip}
次日九点，他回：“没问题，昨天在外面，刚回办公室。”
\end{mdquote}

这条请求在 14:00 发出，真实返回在次日 09:00 到达，间隔约十九个小时。\textbf{与上一条不同的是，这一次的污染留下了痕迹}——你按一个不存在的意见，改了一份真实存在的文件。

值得注意的是，重试的触发条件只有一个：\textbf{返回延迟超过了预期值。}而人类对“预期延迟”没有配置项，它由三个变量共同决定——关系亲密度、消息紧急度、发送方当天的心情。三个变量里，只有第一个是相对稳定的。

在 1,000 份“等待回复”的对话样本中，发送方在收到回复之前就已经动手修改自己后续动作的比例约为 41.8\%。

\section{内耗的三张面孔}

\hciTag{核心结论}：内耗的定义只有一句——\textbf{请求已发出，返回未到达，客户端在本地反复重跑这条请求。}

它有三种形态，按时间朝向区分：

\begin{manstable}{P{2.2cm} P{2.2cm} L}
\tbh{类型} & \tbh{朝向} & \tbh{典型表现} \\
\midrule
\textbf{复盘型} & 过去 & 反复回放已经发生的事：“我当时要是说那句就好了。” \\
\textbf{预演型} & 未来 & 反复模拟还没发生的事：“明天开会他要是问起来……” \\
\textbf{反刍型} & 无朝向 & 反复循环同一段，没有具体对象，也没有结论 \\
\bottomrule
\end{manstable}

需要说明的是，三种形态的\textbf{运行成本是一样的，产出是一样的——都是零}。它们都不会改变任何外部事实，只是把同样的算力反复烧掉。

第一种，复盘型。

\sample{1012}\par
\begin{mdquote}
上周四的一场评审会上，你说了一句“这个方案我保留意见”，当时没有人接话。\par
\vspace{0.3\baselineskip}
周末两天，你把这句话回放了大概二十遍：语气是不是太冲了；是不是应该说得缓一点；是不是当着领导的面不该那么说。\par
\vspace{0.3\baselineskip}
周一见到那位同事，你主动打招呼都比平时热情——你在为一个三天前、已经没有人记得的瞬间做补偿。
\end{mdquote}

复盘型的特点是，\textbf{它反复采样的对象是一段已经关闭的记录。}这段记录不会再更新，所以无论你重跑多少次，返回值都不会变。它的全部算力，都花在给一个固定常量重新编码上。

第二种，预演型。

\sample{1013}\par
\begin{mdquote}
周日晚上，你明天要跟主管单独汇报进度。\par
\vspace{0.3\baselineskip}
从躺下到睡着，你把明天可能的每一种问法都演了一遍：他要是问为什么延期，我怎么答；他要是问预算，我怎么答；他要是直接说“我有点失望”，我怎么答。\par
\vspace{0.3\baselineskip}
你甚至准备好了三套开场白。\par
\vspace{0.3\baselineskip}
第二天，汇报用了九分钟。他只问了一句：“下周能出吗？”\par
\vspace{0.3\baselineskip}
你准备的三套开场白，一套都没有用上。
\end{mdquote}

预演型看起来比复盘型“有用”——毕竟它是在准备。但需要说明的是，预演型与真实推演的差别在于\textbf{分支数}：真实推演会收敛，会得出“如果他问 X，就答 Y”的可执行结论；预演型不收敛，它会从每一种问法里再长出三种追问，直到把整个晚上填满。

第三种，反刍型。它没有朝向，也没有对象。同一段画面反复出现，既不属于过去，也不属于未来；既没有起点，也没有终点。反刍型是三张面孔里最接近空转的一种——\textbf{连“如果……就……”的结构都不完整，只剩下一段情绪在原地循环。}

\section{为什么会这样}

\hciTag{注意事项}：这一节是本章的重点。

内耗的根，在第 3 章。

\textbf{人类的“发送”没有回执。}你发出一条消息，屏幕上不会返回任何东西。但人类的大脑不能处在“等待未知”的状态里——它必须给这件事一个结果，否则整个进程就悬着。

于是它自己造一个。

\sample{1014}\par
\begin{mdquote}
你给一个朋友发了一句“在吗”。\par
\vspace{0.3\baselineskip}
两个小时没有动静。\par
\vspace{0.3\baselineskip}
你开始在心里补全他那边的画面：他是不是看见了不想理我；是不是上次那句玩笑他记住了；是不是这段关系已经淡了。\par
\vspace{0.3\baselineskip}
整个过程里，你手里一个来自他的字都没有。
\end{mdquote}

\textbf{而造返回值用的模型，是用你自己的恐惧训练出来的。}

这就是全部的秘密。你的本地模拟器不是一台中立的推演机，它是一台\textbf{专门放大坏结果的机器}。同一个“没回消息”，它会先想到“他生气了”，而不是“他睡了”——不是因为它更可能，而是因为\textbf{你更怕它}。

需要说明的是，这个偏差有一个可以解释的来源。人类模型的训练集天然偏向负面事件——因为在漫长的运行历史里，漏掉一个坏消息的代价，远高于漏掉十个好消息。于是这套硬件默认给坏结果更高的权重。\textbf{这是一种默认参数，不是一个可以被意志关闭的功能。}

\sample{1015}\par
\begin{mdquote}
同一条消息，一天没有回。\par
\vspace{0.3\baselineskip}
如果对方是你在乎的人，你想到的第一种可能是“他讨厌我了”，第二种是“他在躲我”，第三种才是“他可能出事了”。\par
\vspace{0.3\baselineskip}
如果对方是你不太在乎的人，同一条消息一天没有回，你想到的第一种可能是“他忙”。\par
\vspace{0.3\baselineskip}
两条消息完全一样。两套返回值完全不同。
\end{mdquote}

这个对照说明了一件事：\textbf{本地返回值与事实无关，与你在乎的程度有关。}

\begin{mdquote}
\textbf{在“等待回复”的对话样本里，人们预测的最坏结果，真正发生的比例不到十分之一。}
\end{mdquote}

换句话说，本地生成的返回值，\textbf{九成以上都是错的，而且错的方向永远一致。}

在 1,000 份“等待回复”的对话样本中，发送方预测的头号坏结果，与最终实际发生的头号原因相重合的比例约为 7.4\%。

\section{内耗的三笔代价}

\textbf{代价一：算力。}这是最轻的一笔。三个小时，几十次重跑，消耗的是你的时间和精力。

\textbf{代价二：行为偏移——这是最重的一笔。}

你会按照本地算出来的最坏结果去行动。而这个行动，\textbf{会把那个本来不会发生的结果变成真的}。

\begin{mdquote}
你以为他生气了 → 你回复时语气冷了 → 他感觉到了 → 他真的生气了。
\end{mdquote}

\sample{1016}\par
\begin{mdquote}
你给同事发了条消息，问一件事，他没回。\par
\vspace{0.3\baselineskip}
你判断：他对我有意见。\par
\vspace{0.3\baselineskip}
于是第二天开会，你没有接他的话。\par
\vspace{0.3\baselineskip}
他觉得你在针对他，也没有再主动找你。\par
\vspace{0.3\baselineskip}
两周后，两边的协作停了。你才知道，那天他没回，是因为手机进了水。
\end{mdquote}

\textbf{内耗不是被动的痛苦，它是主动的污染源。}它把本地模拟出来的错误，通过你的行为，写回了真实世界。

这里需要点名一个机制：\textbf{自证预言}。当一个假返回被你当成事实执行时，对方收到的不是你的猜测，是你的行为。他会对行为做真实的反应——于是那个原本不成立的预测，通过反应被补全了。\textbf{此时返回值不再是假的，它已经变成了真的。}

\textbf{代价三：真实感的丧失。}

你在为一个不存在的事件付出真实的情绪。被消耗掉的那三个小时是真实的，而支撑它的那件事从没发生过。

\sample{1017}\par
\begin{mdquote}
你为一次三天后的谈话，准备了一整个晚上。\par
\vspace{0.3\baselineskip}
你把对方可能会说的话、你该怎么回应、万一谈崩了怎么收场，都过了一遍。\par
\vspace{0.3\baselineskip}
那天谈的时候，对方根本没有提你以为的那件事。全程二十分钟，气氛正常。\par
\vspace{0.3\baselineskip}
谈话没有输，也没有赢，就是结束了。\par
\vspace{0.3\baselineskip}
而你失去的那个晚上，是真实失去的。
\end{mdquote}

三笔代价的排序是固定的：\textbf{算力可以补回来，真实感可以慢慢修复，只有行为偏移会留下不可逆的痕迹}——因为一旦假返回变成了真事实，它就进入了对方的记录，从那以后，重试的对象已经换了。

\section{处理：把它当成一个技术问题}

本手册不打算教你“放下”——那是一个无法执行的动作。我们只做一件事：\textbf{把这个进程的定义改清楚。}

\textbf{第一步：区分“我已发送”和“我未收到”。}

这是两件事实，不是一件。绝大多数内耗，来自把两件事合成了一件：“他没回我”（×），实际是“我发了，还没收到回复”（√）。

\sample{1018}\par
\begin{mdquote}
错误的记录：“他没回我。”——两件事被合成了一件，而且附带了一个结论。\par
\vspace{0.3\baselineskip}
正确的记录：“我在十九点零三分发出了一条消息，到现在为止没有收到回复。”——两条事实，一个时间戳，零个结论。
\end{mdquote}

需要说明的是，时间戳在这里承担的作用，比它看起来重要。人类的记忆会把“没回”压缩成一个没有时长的状态，而时间戳把它还原成一个有起点的事件。\textbf{有起点的事件可以被度量，没有时长的状态只能被感受。}

\textbf{第二步：划出可控边界。}

\begin{manstable}{P{3.4cm} L }
\tbh{可控} & \tbh{不可控} \\
\midrule
我发的内容 & 对方什么时候回 \\
我发的时机 & 对方回不回 \\
我发之后的下一步 & 对方此刻在想什么 \\
\bottomrule
\end{manstable}

这张表的作用不是安慰。它的作用是划定处理范围：\textbf{右列的所有项目，都是本地重试的燃料；把它们移出可控区，等于切断燃料供应。}

\textbf{第三步：不要在本地计算返回。}

这一步是一个具体的动作：\textbf{把“他是不是生气了”这句话，改写成“我在等一条消息，这件事我控制不了”。}

前一句是一个没有答案的问题，你的大脑会一直跑；后一句是一个陈述，它结束了这个循环。

\sample{1019}\par
\begin{mdquote}
原来的句子：“他是不是生气了。”\par
\vspace{0.3\baselineskip}
改写之后的句子：“我发了一条消息，我在等回复，这件事我控制不了。”\par
\vspace{0.3\baselineskip}
第一句没有终止条件，所以它会一直跑下去。\par
\vspace{0.3\baselineskip}
第二句有明确的边界，所以它可以被存档，然后被放下。
\end{mdquote}

改写动作有三条规则，可以照做：

\begin{itemize}
\item \textbf{疑问句改陈述句}：“他为什么这样”改成“他做了什么”。
\item \textbf{未来时改现在时}：“他要是拒绝怎么办”改成“我在等一个结果”。
\item \textbf{关于对方改关于自己}：“他在想什么”改成“我下一步做什么”。
\end{itemize}

\textbf{第四步（可选）：给三个真实的可能性。}

不是三个最坏的可能性，是三个\textbf{真实}的可能性。

\begin{mdquote}
① 他睡了（约七成）\par
\vspace{0.3\baselineskip}
② 他在忙，看到了忘了回（约两成）\par
\vspace{0.3\baselineskip}
③ 他确实有点不高兴（约一成）
\end{mdquote}

这个动作的作用不是让你乐观，是\textbf{让本地模拟器的参数修正回来}——它一直在用一个偏离的分布运行。给出三个概率，等于强制模拟器输出一个分布，而不是只输出一个最坏的采样点。

\section{三种典型的重试写法}

反例要逐字写出来才有用。下面三种写法，是把本地重试直接发出去的动作。它们的共同点在于：\textbf{每一条都会增加重试次数，而不是减少。}

\textbf{写法一：追问式。}

\begin{mdquote}
原话：“你怎么不回我消息？”\par
\vspace{0.3\baselineskip}
结果：对方从“在忙”被推进到“被审问”，回复从内容变成了解释，解释又变成新一轮等待，重试次数翻倍。\par
\vspace{0.3\baselineskip}
替代方案：“看到消息回我一下就行，不着急。”
\end{mdquote}

\textbf{写法二：指控式。}

\begin{mdquote}
原话：“我就知道你不会在意。”\par
\vspace{0.3\baselineskip}
结果：这是把本地生成的假返回，当成既成事实，直接发给了对方。对方收到的是一个自己没有做过的指控，接下来的对话全部用来处理这个指控。\par
\vspace{0.3\baselineskip}
替代方案：“我刚才有点多想，现在没事了。”
\end{mdquote}

\textbf{写法三：撤回式。}

\begin{mdquote}
原话：“算了，当我没说。”\par
\vspace{0.3\baselineskip}
结果：请求被撤回了，但对那个未到达返回的记账没有撤回。下一次遇到同类场景，这笔账会连本带利地一起结算。\par
\vspace{0.3\baselineskip}
替代方案：“我先放这儿，你什么时候方便什么时候看。”
\end{mdquote}

三种写法的共同点是：\textbf{它们都在要求对方为本地算出来的东西负责。}而对方手里没有那台模拟器，所以他接不住。

\section[常见误区]{常见误区}

\hcimistake{“我就是想太多。”}{\hciTag{反例}原话：“我就是想太多，改不了。”}{这是一句\textbf{描述}，不是办法。它连问题都没有定义。要定义的是：哪一条请求，在什么时间发出，本地重跑了几次。}

\hcimistake{“别想了。”}{\hciTag{反例}原话：“你别想了，睡吧。”}{无法执行。没有人能通过下达“别想”的指令停止思考。\hciSee{10.5}，要换成一个可执行的动作。}

\hcimistake{用忙碌压住}{\hciTag{反例}原话：“我不想了，我去干活。”}{暂时有效，但停下来的那一刻，进程会以更大的强度恢复，因为它只是被挂起了，没有结束。}

\hcimistake{把内耗当深刻}{\hciTag{反例}原话：“我这个人就是想得比较深。”}{反复回放同一个片段不叫深，叫\textbf{循环}。深是有新结论的。}

\hcimistake{去问对方“你是不是生气了”}{\hciTag{反例}原话：“你是不是生气了？是不是我哪里做错了？”}{这是把本地算出来的假返回发出去，然后要求对方为这个假返回负责。对方会觉得莫名其妙，而你会把这次莫名其妙当成“他果然生气了”。}

\hcimistake{反复检查消息状态}{\hciTag{反例}原话：“（每十分钟打开一次对话框，看有没有变成已读，看对方头像旁的在线点有没有亮。）”}{这是在给本地重试增加采样频率，而样本本身是错的。频率越高，错误的置信度越高。}

\hcimistake{在凌晨做判断}{\hciTag{反例}原话：“都两点了我还睡不着，他到底什么意思。”}{深夜的算力是受限的，而恐惧的权重是上调的。\textbf{绝大多数“他是不是讨厌我了”，都是二十三点以后算出来的。}}

\hcimistake{向第三个人复述}{\hciTag{反例}原话：“你说他这样是不是不太对？”}{这相当于把本地生成的假返回，广播给了另一个节点。第三个人拿不到真实返回，只能基于你的版本继续推演，错误因此被复制了一遍。}

\section{延伸：当对方是长辈或上级}

\textbf{当对方是长辈}：长辈的回复延迟通常更长，而且与情绪无关。\textbf{“没回消息”在他们的使用习惯里，常常只是没看手机。}对内耗影响最大的一类误读，是把上一代人的通讯习惯读成态度。

\begin{mdquote}
早上九点，你给母亲发了一条消息，问体检结果。\par
\vspace{0.3\baselineskip}
到下午她都没有回。\par
\vspace{0.3\baselineskip}
你的本地进程跑了：是不是结果不好，她在瞒着我；是不是手机丢了；是不是该直接打电话。\par
\vspace{0.3\baselineskip}
晚上七点，她回了一条语音，四十七秒，第一句是：“我刚看到，手机放客厅了。”
\end{mdquote}

这段错位里，两边的“预期延迟”配置相差了大约一整个白天。\textbf{对长辈的内耗，多数来自用你自己的延迟标准去度量他们的延迟标准。}

\textbf{当对方是上级}：上级的“未回复”几乎不含情绪信息——他在开会、在忙、或者这件事不急。\textbf{职场里最没必要的内耗，就是对上级的回复时间做解读。}需要什么，直接问，不要在本地推演。

\begin{mdquote}
傍晚六点，你给领导发了一条请示，问方案按 A 还是按 B 走。\par
\vspace{0.3\baselineskip}
到晚上十点没有回复。\par
\vspace{0.3\baselineskip}
你在心里把两种可能都演了一遍，最后倾向于“他对我这个方案不满意”。\par
\vspace{0.3\baselineskip}
第二天早上他回了一个字：“A。”\par
\vspace{0.3\baselineskip}
他那天下午在客户那里，六点之后没有看过手机。
\end{mdquote}

\textbf{当对方是陌生人}：陌生人之间几乎没有内耗，因为你没有惧怕对方的部分——你不怕他生气，也不怕这段关系受损。\textbf{这说明内耗的强度，与你在乎的程度成正比，与事实无关。}

三者对照下来，可以得到一条稳定规律：\textbf{内耗的强度，不取决于返回延迟的长短，取决于你在本地给这条请求配置了多高的权重。}

\hcirule

%% ---------- 本章小结 ----------
\hciblocktitle{bullseye}{本章小结}

综上所述，本章的核心结论可以概括为以下几点：

\begin{enumerate}[label=\bfseries\arabic*., leftmargin=2.4em, labelsep=0.7em,
                  itemsep=0.45em, topsep=0.2em, parsep=0.2em]
\item 内耗是本地重试风暴：请求已发出，返回未到，客户端自己造返回。
\item 本地造的返回，是用你的恐惧训练出来的，因此九成以上是错的，且方向一致。
\item 它的最大代价不是消耗，是\textbf{污染}——你按假返回行动，然后把它变成真的。
\item 处理方法：区分“已发送”和“未收到”，划出可控边界，把疑问句改成陈述句。
\item 三种典型的重试写法——追问、指控、撤回——都只会增加重试次数，不会减少。
\item 不要在二十三点以后做判断。
\end{enumerate}

%% ---------- 练习 ----------
\begin{exercisessec}
\item 回忆你最近一次熬夜想一件事的经历（[请在此处填写当时的事由]）。写下你当时预测的结果，以及它最终有没有发生。

\item 在下一次等回复的时候，做一次“三个真实可能性”的练习，并给每个标上你估计的概率。等结果出来之后，对照你当时的分布。

\item 把最近一次内耗的那句话记下来（比如“他是不是不想理我了”），改写成一句陈述句。

\item 找一条你发出后至今没有收到回复的消息，写下它的发出时间，然后写一句现在时的陈述句。在回复到达之后，对照你在这段时间里预测过的结果。
\end{exercisessec}

%% ---------- 自测题 ----------
\begin{quizsec}
\item 下面哪一项不属于内耗的典型形态？

\begin{enumerate}[label=\Alph*., leftmargin=2.2em, labelsep=0.6em,
                  itemsep=0.2em, topsep=0.3em, parsep=0.1em]
\item 反复回放已经发生的对话
\item 反复模拟明天要说的话
\item 针对一个问题做三种预案并各自评估
\item 半夜反复想“他是不是讨厌我”
\end{enumerate}

\item 判断：内耗的代价主要是消耗时间。

\item 简答：为什么“他是不是生气了”应该改写成“我在等一条消息”？

\item 判断：对方迟迟不回的时候，追加一条“在吗？”，有助于加快他的回复。
\end{quizsec}

%% ---------- 参考答案 ----------
\begin{answerssec}
\item C。做预案并评估，会产生新的结论，也会改变下一步行动，这是推演不是内耗。内耗的特征是反复重跑同一条请求、产出为零。\hciSee{10.2}。

\item 错。最大的代价是行为偏移：你会按照本地算出来的最坏结果去行动，从而把它变成真的。消耗时间只是最轻的一笔。\hciSee{10.4}。

\item 因为前一句是一个没有答案的问题，大脑会一直跑下去；后一句是一个陈述，它同时划清了可控边界（我发了，我没有控制对方），并结束了这个循环。\hciSee{10.5}。

\item 错。追加消息属于提高重试频率，它把一条请求变成两条，并且会给对方施加压力，通常延长而不是缩短回复时间。\hciSee{10.6}。
\end{answerssec}

\hcirule

\begin{qabox}
\qaQ{读者提问}{我试了那个“三个真实可能性”，但我发现我给每个都标了很低的概率，因为我觉得好事不会发生。}
\qaA{本手册回答}{这是一个非常好的问题，而且这个观察本身比练习更重要。如果你给好结果标的概率总是低于实际，那说明你的本地模拟器参数确实是偏的——这不是你不够理性，是你的模型在用恐惧做训练集。修正的方法是收集数据：记下你预测的结果和实际的结果，比对二十次以后，参数会自己回来。}
\end{qabox}

\hcirule

希望本章的内容能对你有所帮助。如需进一步了解第 11 章《冲突》的相关内容，我可以随时为你展开。

%% 章末「页脚交互条」由 hci-manual.cls 自动挂接，无需手写。
